% Generated by Listas/make_statements.py from pratica-lecture-07.md, item 17. Do not edit.
\textbf{Source.} Tadelis (2013), \emph{Game Theory: An Introduction}, Exercise 13.3 ``Reserve Prices (Easier)'', pp.~285--286 {[}PDF pp.~304--305{]}. Reproduced verbatim. \textbf{Answer in the source:} no (none in the book; Princeton UP's website has about half, not checked).



\begin{quote}

\textbf{Reserve Prices (Easier):} Consider a seller who must sell a single private value good. There are two potential buyers, each with a valuation that can take on one of three values, \(\theta_i \in \{0, 1, 2\}\), each value occurring with an equal probability of \(\frac{1}{3}\). The players' values are independently drawn. The seller will offer the good using a second-price sealed-bid auction, but he can set a ``reserve price'' of \(r \ge 0\) that modifies the rules of the auction as follows. If both bids are below \(r\) then neither bidder obtains the good and it is destroyed. If both bids are at or above \(r\) then the regular auction rules prevail. If only one bid is at or above \(r\) then that bidder obtains the good and pays \(r\) to the seller.



\begin{enumerate}

\def\labelenumi{\alph{enumi}.}

\item

  Is it still a weakly dominant strategy for each player to bid his valuation when \(r > 0\)?

\item

  What is the expected revenue of the seller when \(r = 0\) (no reserve price)?

\item

  What is the expected revenue of the seller when \(r = 1\)?

\item

  What explains the difference between your answers to (b) and (c)?

\item

  What is the optimal reserve price \(r\) for the seller, and what can you conclude about the value of reserve prices?

\end{enumerate}

\end{quote}

